\section{Two estimates for the stationary configuration}
\label{sec:stationary-estimates}

We use the underlying equality and the collapsed event law of
Proposition~\ref{prop:stationary-data}. They have different purposes:
plank bounds always count inherited weights of original active indices,
whereas probabilities below are probabilities of indexed events. In
particular, restricting an event label does not restrict the reference
index mass in a plank test.

We prove two estimates. Lemma~\ref{lem:angular-frostman} bounds the
concentration of normal labels along an index and gives $\beta<1$.
Lemma~\ref{lem:conditional-X} bounds the additional normal-position
entropy inside a matrix cell after conditioning on a finer event label.
The first estimate follows by turning angular concentration into a
configuration with excessive temporal deficit. For the second, we
compare the mass of a queried packet with the original active-index
mass in its matrix cell, and then convert that comparison to entropy.

Throughout this section put
\[
 P=P_{p,q},\qquad x=d-1-h,\qquad y=1+h,
 \qquad x+y=d,\qquad y-h=1.
\]
The parameters $0<\tau,\ell<1$, $\tau+\ell=1$, $v>0$,
$\beta=v/\ell\le1$, and $0<w\le1$ are fixed. Thus
$\beta>v$, including before we exclude $\beta=1$.

Recall the four matrix widths $L_X(c),L_Y(c),L_U(c),L_V(c)$ from
\eqref{eq:estimate-widths}. For a coarse event label
$e_c=(t_c,\vartheta_c)$, evaluate the four forms at its chosen center.
For a real offset $a$, let $Q(c,a)$ be their joint matrix-cell label
at those widths multiplied by $N^{-a}$. The grids are dyadic, with the usual
bounded enlargements at rounded scales. Conditional on $e_c$, this is
a partition of the matrix indices; repetitions at the same matrix
retain their individual weights. All limits first fix the scale
parameters and any strict exponent margin, and then let $N\to\infty$.
Finite lists of such tests precede the countable diagonals.

\subsection{Mass accounting for the selections}

Both proofs select events and then fill the occupied time blocks they
meet. We must retain the packet mass and control the change in the time
profile under the same selection. The following form of finite cleaning
does this even when the selection has a fixed power loss. Its essential
distinction is between original packet masses, against which that loss
is measured, and the restricted graph on which new probability labels
are computed.

\begin{lemma}[Cleaning registered events]
\label{lem:estimate-cleaning}
Let a finite weighted event graph have total weight $W$, and let a
restriction retain weight at least $N^{-E+o(1)}W$. Fix finitely many
partitions of the graph. The cells may include packet assignments,
index--time blocks, and joint labels. After a further restriction of
subpower cost, every surviving
cell of each partition retains at least $N^{-E-o(1)}$ of its original
weight. The same assertion holds after splitting each original cell
among subpower many registered charts, with that number absorbed in
the error.

Suppose, in addition, that the original time trees have a common
profile $F$ on a log interval $[0,\Lambda]$ in $\log N$ units, and each
original index and starting interval is routed through only subpower
many charts. If the restricted trees are made uniform, with profile
$G$, then
\begin{equation}
 G(t)-G(s)\ge F(t)-F(s),\qquad
 0\le G(\Lambda)-F(\Lambda)\le E,
 \quad 0\le s<t\le\Lambda.
\label{eq:estimate-branch-loss}
\end{equation}
In particular an extensive restriction preserves the profile, and
\begin{equation}
 0\le \int_0^\Lambda P(G')-\int_0^\Lambda P(F')\le pE.
\label{eq:estimate-cost-loss}
\end{equation}
These assertions concern inherited index weights.
\end{lemma}

\begin{proof}
We distinguish two reference graphs in the cleaning. For the new
homogenization tests use the finite label construction of
Lemma~\ref{lem:regularization}: these are its within-index integer
counts and conditional-probability tests with polynomially many
children. Perform its negligible conditional-probability tail
deletions before computing the labels. Arbitrary-cardinality original
packet partitions are retained only as cleaning tests; no absolute
weight label is imposed on their cells. Let $W$ be the original event
weight and let $R$ be the restricted, routed graph after those tail
deletions and before the finite labels are selected. Write $M_0=|R|$;
thus $M_0\ge N^{-E-o(1)}W$. Compute the prescribed rounded labels
on $R$, including their joint and domain tests for conditional counts,
and retain one common finite list of labels. All graph masses denoted
by $|\cdot|$ here and in the registration rule below are weighted
event masses. Denote this label
class by $A$, and write $M=|A|$. The number of label lists is subpower,
so $M\ge N^{-o(1)}M_0\ge N^{-E-o(1)}W$. Counts within index fibers
are event counts, with the inherited index weight cancelling in their
ratios. No bound on the number of distinct weighted indices is required.

There are two collections of reference masses. Original packet,
index--time, and chart tests use their masses in the original graph;
the sum of these references is at most $L_{\rm old}W$.
New regularization tests use their masses in $R$, the graph on which
their labels were computed; their reference sum is at most
$L_{\rm new}M_0$. Here $L_{\rm old},L_{\rm new}=N^{o(1)}$
include the finite number of partitions and any subpower chart copies.
In particular, one does not substitute the possibly much smaller
post-label cell mass for an $R$-reference without recomputing its label.

Choose $\epsilon_N\downarrow0$ slowly enough to dominate the label
loss $\log_N(M_0/M)$, the error in $M\ge N^{-E-o(1)}W$, and the
$\log_N L_{\rm old}$ and $\log_N L_{\rm new}$, with margins tending to
infinity after multiplication by $\log N$. Starting from $A$, delete
all events in a deficient original-reference cell when its current
mass is less than $N^{-E-\epsilon_N}$ times that reference. For a
new regularization cell or domain, use instead $N^{-\epsilon_N}$
times its $R$-reference. Continue either type of deletion until no
deficient cell remains. A cell is charged only once, because it is
empty after deletion. Even with cascades, the two total charges are
at most
\[
 L_{\rm old}N^{-E-\epsilon_N}W=o(M),\qquad
 L_{\rm new}N^{-\epsilon_N}M_0=o(M).
\]
The final graph therefore has mass $(1-o(1))M$. Every surviving old
cell retains at least $N^{-E-o(1)}$ of its original mass. Every
surviving new cell retains $N^{-o(1)}$ of its $R$-mass and is a
subset of that cell, so the rounded exponents selected from $R$
remain valid, including conditional exponents obtained by ratios of
joint and domain masses. This proves the mass assertion: the new labels
incur only a subpower loss in addition to the original loss $E$. The
finite-test diagonal allows further prescribed tests.

For the profile assertion fix a finite mesh of time scales and
uniformize the branching counts on its gaps. A subtree cannot have
more children on any gap than the original tree. Since the branching
exponent on $[s,t]$ is $(t-s)-(F(t)-F(s))$, this proves the first
inequality in \eqref{eq:estimate-branch-loss}. The total available
index weight after routing is at most $N^{o(1)}$ times its original
value. Comparing this with the retained weighted event count shows
that the number of events per surviving uniform index is at least
$N^{-E-o(1)}$ times the original number. This proves the endpoint
inequality. Refine the fixed mesh only after these comparisons hold.
The Lipschitz profile limits give \eqref{eq:estimate-branch-loss} on
every interval, hence $G'\ge F'$ almost everywhere and
$\int(G'-F')\le E$. The function $P$ is increasing and has Lipschitz
constant at most $p$, which proves \eqref{eq:estimate-cost-loss}.
\end{proof}

Here is the registration rule used below. Before any restriction, an
occupied time block $J$ on an assigned index has a common number
$m_J=N^{\kappa+o(1)}$ of finer events, uniformly in the blocks under
consideration. Given a selected set of finer events, first register
every occupied token $(i,J)$ that it hits, giving that token weight
$\omega_i$, and complete its block. Any preassigned chart or other
label on which the predicate depends remains part of its token;
completion never changes that label. Subpower many such tags enlarge
the original reference totals by only a subpower factor. If $T_0$ is the original token
weight and $T$ the weight of the registered tokens, then
\[
 |\text{selected finer events}|
 \le N^{\kappa+o(1)}T,\qquad
 |\text{original finer events}|=N^{\kappa+o(1)}T_0.
\]
Consequently a restriction of finer-event weight $N^{-E-o(1)}$
registers at least $N^{-E-o(1)}T_0$ token weight. Completing these
tokens precedes chart splitting and incurs no further loss. Route
whole tokens through their geometric charts; each original index and
starting block has only subpower many such charts. If it is necessary
to prune chart portions, compare them with their masses in this
restricted, completed token graph, rather than with the original
unrestricted block's finer-event mass. Portions below
$N^{-\epsilon_N}$ of their restricted unsplit reference have total
weight $o(T)$ after choosing $\epsilon_N$ slowly. A predicate constant
on $(i,J)$ is unchanged by completion or later filling. The original
packet references are still the unrestricted assigned token masses,
and keep the $N^{-E-o(1)}$ threshold in the preceding two-budget
cleaning. Thus this registration retains both the original packet
lower bounds and the single event-retention loss $E$.

\subsection{Angular nonconcentration}

\begin{lemma}[Angular Frostman bound]
\label{lem:angular-frostman}
There is $\gamma>0$, depending only on the fixed stationary parameters,
with the following property. For every fixed $0<c<b\le1$, outside a
set of index--$I_c$ fibers of event probability tending to zero, every
angular interval $J$ of length $\Theta_b$ satisfies
\begin{equation}
 \frac{\#\{e\text{ on }i:t(e)\in I_c,
                         \ \vartheta(e)\in J\}}
      {\#\{e\text{ on }i:t(e)\in I_c\}}
 \le N^{-\gamma\tau(b-c)+o(1)}.
\label{eq:angular-frostman}
\end{equation}
It is enough to test a fixed grid with bounded enlargements. The bound
also holds conditional on typical nondeficient $e_c$ bins in the fiber,
and after extensive regular refinements. In particular
\begin{equation}
 0<\beta<1,\qquad s_0=1-\beta>0.
\label{eq:positive-time-dimension}
\end{equation}
\end{lemma}

\begin{proof}
Angular concentration would produce a configuration whose normalized
temporal deficit is at least $\beta$. Since $\beta>v=-h_p$, decreasing
$p$ makes the critical upper bound for that configuration strictly
smaller than its packet lower bound. We quantify how close the two
bounds remain after selecting a narrow angular bin: a selection loss
$E$ decreases the packet mass by at most $E$ in exponent and increases
its temporal cost by at most $pE$.

Write $g=b-c>0$ and
\[
 S=n_bN^{-\tau b},\qquad K'=N^{b-\ell c},
 \qquad\theta=\Theta_b=N^{-\tau b}.
\]
First pass to the isotropic outer parents associated to $B_b$, as in
Lemma~\ref{lem:nested-cuts}. Their effective resolution count is
$N^b$. The assigned $B_b$ plates have time and long spatial length
$N^{\tau b}$ in these cell units, unit thin width, and mass at least
$SN^{\tau b-o(1)}$. On the longer interval $I_c$, which has $K'$ bins
in this model, the local bound of Proposition~\ref{prop:local-bound}
pulls back to
\begin{equation}
 n(\text{plank of widths }a_1,b_1)
 \le N^{o(1)}
 S N^{-\ell gP(\beta)}(K')^{-h}a_1^xb_1^y.
\label{eq:angular-parent-bound}
\end{equation}
For example, the isotropic physical cell width is
$N^{w(1-b)}$ in original cell units. Substitution in the original
local bound gives the coefficient
\[
 mN^{-\ell(1-c)P(\beta)}N^{-h(1-\ell c)}
       N^{dw(1-b)}
 = S N^{-\ell gP(\beta)}(K')^{-h}
\]
in exponent, using $m\eqexp\Delta_hN^{B_*}$ and
\eqref{eq:stationary-packet-mass}. Thus
\eqref{eq:angular-parent-bound} counts original active indices; it is
not a bound only for indices occupying a selected angular bin.

Suppose that fibers violating \eqref{eq:angular-frostman} account for
an extensive event piece. On each such index and $I_c$ choose one
violating bin and keep its events. This costs at most
\[
 E=\gamma\tau g
\]
in log probability, apart from a vanishing error. First register and complete every occupied $I_b$ token hit by this
selection, before any chart-portion cutoff. The preceding registration
rule gives token weight at least $N^{-E-o(1)}$ of the original token
graph. The normal of its assigned $B_b$ packet is within
$\Theta_bN^{o(1)}$ of the selected bin, since the packet contains the
event that registered the token and its normal label is constant on
the token. Completing the block does not change its packet owner.
Route these whole tokens through the outer parents and bin choices.
Each packet has only subpower many relevant choices, and each index
on $I_c$ has only subpower many charts. Any deficient chart portions
are pruned relative to this restricted token graph. Now apply the
two-budget cleaning of Lemma~\ref{lem:estimate-cleaning}: original
packet and time-tree references use the loss $E$, while the new
homogenization references, computed after registration and routing,
use only a subpower loss. The final token graph still has weight
$N^{-E-o(1)}$ of the original one, and each surviving packet has raw
assigned mass at least
\begin{equation}
 SN^{\tau b-E-o(1)}.
\label{eq:angular-selected-packet}
\end{equation}
The uniform token trees are subtrees of the original occupied
$I_b$ trees below $I_c$, with subpower many routed copies. Thus their
endpoint deficit loss is at most $E$ and their cost increase is at
most $pE$. Packet loss $E$ and cost increase $pE$ are the two
consequences of this same retained graph; the selection is not
performed a second time.

Apply the geometric construction in the proof of
Lemma~\ref{lem:anisotropic} on $I_c$, with angular width $\theta$.
Here we use only its bounded-chart routing, rectangle pullback, and
cell-count calculation. That calculation is linear in the packet
mass, which here is the reduced value $SN^{\tau b-E-o(1)}$ from
\eqref{eq:angular-selected-packet}. We do not use that lemma's
profile-preservation conclusion for an extensive graph: the endpoint
loss $E$ and cost increase $pE$ were proved by
Lemma~\ref{lem:estimate-cleaning} and are used explicitly in
\eqref{eq:angular-cost} below. The resulting normalization has resolution
\begin{equation}
 N_{\mathrm{new}}=K'\theta=N^{\ell g}.
\label{eq:angular-new-resolution}
\end{equation}
The selected plates become bounded-cell packets: their old time and
long spatial length multiplied by $\theta$ is one, and their
orientation error is $N^{o(1)}\theta$. Filling their occupied $I_b$
blocks before this normalization therefore gives aggregate raw
multiplicity at least the quantity in
\eqref{eq:angular-selected-packet}. The full-plank parameter of the
new models is at most
\begin{equation}
 \Delta_{\mathrm{new},h}
 \le N^{o(1)}S N^{-\ell gP(\beta)}(K')^{-h}\theta^{-y}.
\label{eq:angular-new-clustering}
\end{equation}
Here the factor $\theta^{-y}$ follows by pulling back a rectangle:
its area expands by $\theta^{-1}$ and its largest width by at most
$\theta^{-1}$, and $a^xb^y=(ab)^xb^{y-x}$ with $0\le x\le y$.

We next quantify the time cost. Prior to selection the occupied
$I_b$ blocks below $I_c$ have constant deficit slope $\beta$ over a
log interval of length $\ell g$. Filling and then normalizing removes
only the flat range below them. The restricted uniform time tree
cannot gain branching on any fixed gap. Its normalized deficit
$\alpha_{\mathrm{new}}$ is therefore at least $\beta$. By
\eqref{eq:estimate-cost-loss}, its cost, expressed in $\log N$ units,
is at most
\begin{equation}
 \ell gP(\beta)+pE+o(1).
\label{eq:angular-cost}
\end{equation}
The comparison uses the whole routed ensemble: the number of chart
copies per original index and $I_c$ is subpower, so the total
branching loss is at most $E$, as in the proof of
Lemma~\ref{lem:estimate-cleaning}.

Without the controlled losses, the sharp upper bound from
Definition~\ref{def:critical} and
\eqref{eq:angular-new-clustering} matches $SN^{\tau b}$ exactly.
Indeed its exponent relative to $S$ is
\begin{align*}
 -\ell gP(\beta)-h(b-\ell c)+y\tau b
       +h\ell g+\ell gP(\beta)
 &=\tau b.
\end{align*}
Thus \eqref{eq:angular-selected-packet} and
\eqref{eq:angular-cost} place the new models within
$(p+1)E+o(1)$ of their sharp upper bound in $\log N$ units.

Choose a small fixed $t_*>0$ with $p-t_*>2$ and
\begin{equation}
 0\le h(p-t_*,z)-h(p,z)
 \le \frac{v+\beta}{2}\,t_*.
\label{eq:angular-p-perturbation}
\end{equation}
This is possible by differentiability and $\beta>v$.
Increasing $h$ at fixed $d$ cannot increase the plank parameter,
because its value on a plank is multiplied by
$(a_1/b_1)^{h'-h}\le1$. Also
$P_{p-t_*,q}(e)=P_{p,q}(e)-t_*e$. Applying the critical bound at
$(p-t_*,z)$ to these same models therefore saves at least
\[
 N_{\mathrm{new}}^{(\beta-v)t_*/2}
\]
relative to the old sharp upper estimate. Choose once and for all
\begin{equation}
 0<\gamma<\min\left\{
 \frac12,\frac{\ell(\beta-v)t_*}{8(p+1)\tau}
 \right\}.
\label{eq:angular-gamma-choice}
\end{equation}
The saving then strictly exceeds the error $(p+1)E$. This contradicts
the aggregate lower multiplicity. Uniformity over bounded parent
charts follows from the uniform all-configuration critical bound;
otherwise a parent subsequence would violate that bound.

It follows that the violating fibers are not an extensive piece.
In particular their probability tends to zero. This proves
\eqref{eq:angular-frostman} for each fixed gap. The choice
\eqref{eq:angular-gamma-choice} is independent of $b-c$: both the
selection loss and the strict saving were proportional to that gap
before normalization. Fix finitely many gaps first and absorb their
errors; the diagonal convention then gives all required scale tests.

Within an index and $I_c$ the normals occupy only subpower many
$e_c$ bins. Delete bins retaining less than a subpower fraction of the
fiber, by the same deficiency argument. Dividing by their remaining
mass changes \eqref{eq:angular-frostman} by only a subpower factor.
An extensive subsequent regular refinement has the same property by
Lemma~\ref{lem:estimate-cleaning}.

Finally suppose $\beta=1$. For any fixed $0<c<1$, an occupied
index--$I_c$ fiber has only $N^{o(1)}$ occupied $I_1$ blocks by
\eqref{eq:stationary-time-branching}, and there is one collapsed
event per such block. One angular bin of width $\Theta_1$ therefore
has at least $N^{-o(1)}$ of its event mass. This contradicts
\eqref{eq:angular-frostman} at $b=1$, since $\gamma\tau(1-c)>0$.
Hence $\beta<1$; positivity follows from $v>0$.
\end{proof}

\subsection{Conditional concentration in the normal position}

We next compare a queried packet with the original active-index mass
in its matrix cell. This reference mass includes indices whose shadings
do not carry the query's finer event label. A lower bound for such a
reference mass can nevertheless bound the query's entropy: entropy is
at most cross-entropy against any positive reference probability. The
proof first establishes the geometric mass bound and then applies this
inequality to the normal-position observation.

For a finite random variable write
$\mathrm H_N(Z\mid L)=\mathrm H(Z\mid L)/\log N$, using natural
logarithms. If $L=e_{c+\sigma}$, let $X_L$ be the normal-position
form evaluated at the center of that finer label. Let
\[
 [X_L]_{c,a+\Delta}
\]
denote its grid bin at width $L_X(c)N^{-a-\Delta}$. Its grid may be
translated by a quantity determined by $Q(c,a)$ and $L$; changing
between such grids changes the conditional entropy by $O(1)$.

\begin{lemma}[Conditional $X$ estimate]
\label{lem:conditional-X}
Fix $0<c<1$, $0<a<w(1-c)$, and $\sigma>0$ with
$c+\sigma\le1$. For all sufficiently small fixed $\Delta>0$,
\begin{equation}
 \limsup_{N\to\infty}
 \mathrm H_N\bigl([X_{e_{c+\sigma}}]_{c,a+\Delta}
                  \mid Q(c,a),e_{c+\sigma}\bigr)
 \le x\Delta.
\label{eq:conditional-X-bound}
\end{equation}
For example it suffices to take
\begin{equation}
 0<\Delta<\frac14\min\left\{
 a,\ w(1-c)-a,\ \frac{\tau a}{w},\
 \min(\tau,\ell)\sigma
 \right\}.
\label{eq:conditional-X-admissible}
\end{equation}
The assertion holds for the event law on any extensive regular
refinement, with the original active-index masses still available as
reference masses. The limit in $N$ precedes limits in $\Delta$ or
$\sigma$.
\end{lemma}

\begin{proof}
The geometric argument uses coarse matrix cells within registered
charts. We will replace each cell by its center, weighted by the full
original active-index mass in that cell, and regard the selected
original indices as its children. After selecting a common ratio of child mass
to representative mass, their normalized event laws are comparable.
The critical inequality bounds the representatives' multiplicity,
whereas packet saturation gives a lower bound for the children's mass
in a typical spatial cell. These bounds force a typical packet to
contain a sufficiently large fraction of its queried cell's reference
mass. We first construct the charts in which this comparison is made.

Put
\begin{equation}
 g=b-c=\frac{a+\Delta}{w},\qquad
 D_0=N^g,\qquad D=N^\Delta,\qquad
 Z'=N^{\tau g},\qquad S'=\frac{n_b}{Z'}.
\label{eq:X-test-parameters}
\end{equation}
The conditions \eqref{eq:conditional-X-admissible} imply
$c<b<1$ and $\Delta<\tau g$, so a positive flat range remains in
the time filling used below.

\paragraph{Chart geometry.}
Fix $e_c$ and its interval $I_c$. In its frozen frame normalize time
by $R_c$, spatial normal and transverse position by
$R_c\Theta_c,R_c$, and velocity by the corresponding spatial units
divided by $R_c$. In these coordinates a $Q(c,a)$ cell has isotropic
matrix diameter, up to dimension constants,
\[
 N^{-\rho(1-c)-a}.
\]
Take isotropic matrix parents of diameter $N^{-\rho(1-b)}$, subtract
their centers, and divide the spatial and velocity coordinates by
that diameter. The resulting charts are bounded and have resolution
$D_0^{-1}$. In each such parent, the diameter of a cell of
$Q(c,a)$ is
\begin{equation}
 N^{-\rho(b-c)-a}=N^{-g+\Delta}=D/D_0.
\label{eq:X-cell-diameter}
\end{equation}
We used $wg=a+\Delta$ in the equality. One may choose compatible
nested grids; if a prescribed cell meets a parent boundary, splitting
it among boundedly many parents adds only $O(1)$ entropy at the end.
Hereafter $C$ denotes the resulting cell within its registered chart,
including this subdivision when needed.

A $B_b$ packet fits in subpower many of these charts. Its used
normals lie in $\Theta_bN^{o(1)}$, hence in subpower many $e_c$
bins, and its four width bounds extrapolated from $I_b$ to $I_c$
have matrix diameter at most $N^{-\rho(1-b)+o(1)}$ in the
$I_c$ normalization. Use the packet's own normal and its own time center for the following
four widths after parent normalization. Their negative log exponents
are
\begin{equation}
 (g,\ell g,\tau g,0).
\label{eq:X-normalized-packet-widths}
\end{equation}
In the normalized $I_c$ chart the change of normal is
$(\vartheta_P-\vartheta_c)/\Theta_c=O(1)$, and the change of time
center is $(t_P-t_c)/R_c=O(1)$. These are bounded shears. They need
not preserve the four coordinate widths in the frozen $e_c$ frame;
they preserve the following geometric bounds, up to fixed factors.
At resolution $D_0$ the packet is a plate of thin spatial width one
cell, long spatial width $Z'$, and time length $Z'$ bins, with raw
assigned mass at least $S'Z'N^{-o(1)}$.
Each index and $I_c$ is in only subpower many registered charts.
After chart splitting, Lemma~\ref{lem:estimate-cleaning} preserves
these packet masses on typical occupied pieces.

\paragraph{The reference-mass claim.}
For each cell $C$ in a registered chart define
\begin{equation}
 n_C=\sum_{\substack{i\text{ active on the original }I_c\\M_i\in C}}
                   \omega_i.
\label{eq:X-reference-mass}
\end{equation}
For a packet $P=B_b$ on an $I_b$ block, write $n(C\cap P)$ for
the inherited mass of its assigned active indices that lie in $C$.
We will prove that, outside a set of query events of probability
tending to zero,
\begin{equation}
 n(C\cap P)\ge n_C D^{-x-o(1)}.
\label{eq:X-reference-packet}
\end{equation}
The claim also holds if the query events were first restricted to an
extensive regular piece, with the same full original active mass in
\eqref{eq:X-reference-mass}. This reference mass belongs to the chart cell
$C$ just defined.

The critical comparison uses the following inherited bound. Let
$\Delta_0$ denote the full-time plank parameter in
one of these bounded normalized charts, with exponents $x,y$. Its
physical normal cell width in original cell units is $N^{w(1-b)}$;
the transverse cell width is
$\Theta_c^{-1}N^{w(1-b)}$. Pulling back an arbitrarily oriented
rectangle multiplies its area by
$N^{2w(1-b)}\Theta_c^{-1}$ and its largest width by at most
$N^{w(1-b)}\Theta_c^{-1}$. Since $0\le x\le y$, the original
local estimate gives
\begin{align}
 \Delta_0
 &\le N^{o(1)}mN^{-\ell(1-c)P(\beta)}
       N^{-h(1-\ell c)}N^{dw(1-b)}\Theta_c^{-y}
       \notag\\
 &\eqexp S'D_0^{-B_*}.
\label{eq:X-chart-clustering}
\end{align}
For the last equality, its exponent relative to $\Delta_h$ is
$cB_*+\tau c+dw(1-b)$; the same exponent results from
$S'D_0^{-B_*}$ using \eqref{eq:stationary-packet-mass}.
All these bounds count original indices active on $I_c$, whether or
not those indices carry the particular normal label used to name
the chart. An upper bound for a subcollection is therefore inherited
without changing the reference weights.

\paragraph{Registering the failed queries.}
Suppose the reference-mass claim fails. After a subsequence, for some
fixed $\epsilon>0$
the events satisfying
\begin{equation}
 n(C\cap P)<n_CD^{-x}N^{-\epsilon}
\label{eq:X-bad-predicate}
\end{equation}
form an extensive piece. Register tuples $(i,\mathfrak c,J)$,
where $\mathfrak c$ specifies the chart, including $e_c,I_c$ and
the parent, and $J=I_b$ is occupied. Give each tuple weight
$\omega_i$. The cell $C$ is fixed by $(i,\mathfrak c)$, and its
packet owner is fixed by $(i,\mathfrak c,J)$, by
Proposition~\ref{prop:stationary-data}(iv). Consequently
\eqref{eq:X-bad-predicate} is constant on the tuple. This constancy
is the reason that replacing its finer events by a token and later
filling $J$ does not change the failed test.

Before selection each occupied $J$ has a common number of original
events in exponent. Register every tuple $(i,\mathfrak c,J)$ hit by
a bad query and complete its $J$ block in that fixed chart, before
chart-portion pruning. The chart component is retained throughout:
the failed predicate depends on it. Because the predicate is constant
on this tuple, completion preserves it. The
bad tuples consequently have extensive inherited weight in the
original token graph. Any chart cutoff is taken relative to the
restricted completed token graph; there are only subpower many
chart portions per index and $I_c$. Clean the packet partitions,
using the original token masses for packet references and the graph
where labels were computed for new regularization references.
The registered assigned mass of each surviving packet in a chart
is then at least $S'Z'N^{-o(1)}$. The assigned index sets are
disjoint for a fixed chart and $J$.

\paragraph{Frequency selection.}
We next make the selected child mass comparable to the reference mass
by a common factor. This will allow a multiplicity estimate sampled
under the representative law to be used for the child events.
On the retained tuple graph define
\[
 \mathsf w(C,J)=
   \sum_{\substack{(i,\mathfrak c,J)\text{ retained}\\M_i\in C}}
         \omega_i.
\]
This is at most $n_C$. Select one common exponent of
$\mathsf w(C,J)/n_C$ across the ensemble. Here is the finite
justification even when the original weights are arbitrarily
unequal. There are polynomially many charts, cells, and time blocks,
and each $n_C\le n$. If their total number is at most $N^A$,
ratios below $N^{-A-3}$ contribute at most $N^{-3}n$ to the tuple
weight and may be discarded. The extensive original token graph has
weight at least $nN^{-o(1)}$: every original index had a common total
event count, and the restriction cannot create more events per
index. Partition the remaining logarithmic ratios in $[0,A+3]$ on
a mesh tending slowly to zero. There are subpower many levels, so
one retains an extensive level. After further cleaning the
$(C,J)$ cells, using their masses at this level as references,
there is an exponent $r\ge0$ such that
\begin{equation}
 \frac{\mathsf w(C,J)}{n_C}=N^{-r+o(1)}
\label{eq:X-frequency}
\end{equation}
uniformly on all surviving pairs. It is the selected
$\mathsf w(C,J)$, not $n_C$, that is used as the reference for this
deficiency deletion. Therefore the sum of reference masses in this
cleaning partition is the tuple weight, not a possibly much larger
sum of $n_C$'s. Include the packet cells and all prescribed time-tree
tests in the same deletion process. The packet lower masses, the
ratio \eqref{eq:X-frequency}, and the inherited uniform child
branching then hold simultaneously.

\paragraph{Representatives and their event law.}
Fill every registered tuple over its whole interval $J$ and rebin
time at the coarser resolution
\[
 (D_0/D)^{-1}.
\]
Each $J$ contains $Z'/D$ such bins, up to the harmless dyadic
rounding, and this number has positive exponent. The operation
drops only part of the full-branching range grafted below $J$.
For each occupied $C$ put one representative matrix at its center,
of weight $n_C$. Occupy that representative at every coarse time
bin belonging to a registered pair $(C,J)$. Charts are counted
separately. A child tuple at a coarse bin is assigned to the spatial
cell hit by its representative at that bin center; this is the
hashing used in the argument.

On each representative event the total weight of its children is
$\mathsf w(C,J)$. By \eqref{eq:X-frequency}, the ratio of that
weight to the representative weight is $N^{-r+o(1)}$, uniformly.
The same common factor compares the total masses. Thus the two
normalized event laws are mutually dominated up to subpower
factors. In particular, an extensive set of representative events
pulls back to an extensive set of interpolated child events. This
statement does not identify the individual index weights $n_C$
with the child weights.

\paragraph{The representative multiplicity bound.}
We apply the critical inequality to this representative configuration.
Its plank parameter is inherited from the reference indices; its
temporal cost will be bounded using the selected children.
The full-plank parameter for this representative model is at most
\begin{equation}
 \Delta_{\mathrm{rep},h}
 \le N^{o(1)} S'D_0^{-B_*}D^d.
\label{eq:X-representative-clustering}
\end{equation}
Indeed, by \eqref{eq:X-cell-diameter}, every original line counted
in $n_C$ stays within a bounded coarse cell of the representative
throughout the whole normalized time interval. A coarse plank
containing the representative therefore contains all these lines
after bounded enlargement. In the old $D_0$ cell units its widths
are at most constant multiples of $Da_1,Db_1$. The cells $C$
partition the original active indices within a chart, so summing
their weights does not duplicate any index in that chart. Apply
\eqref{eq:X-chart-clustering} and use $x+y=d$ to obtain
\eqref{eq:X-representative-clustering}.

We claim that the sum of the occupied representative weights in a
typical hashed space--time cell is at most
\begin{equation}
 S'D^{d-h}N^{o(1)}.
\label{eq:X-representative-multiplicity}
\end{equation}
If cells with a strict excess carried extensive representative
incidence weight, restrict to those cells. Their aggregate support
is at most their weighted incidence count divided by the proposed
excess threshold. Make the representative time profiles uniform
over the ensemble, retaining an extensive piece.

It remains to bound the cost of this profile. Up to the
creation of representatives, all selections (the bad predicate,
chart routing, frequency-level selection, and their regularizations)
are extensive restrictions of inherited-weight child tokens. The
reference mass $n_C$ used to define the frequency ratio never becomes
a child index weight. Filling each registered $J$ adds the same
$Z'/D$ coarse events per token. It follows from
Lemma~\ref{lem:estimate-cleaning}, on every tested gap, that the
interpolated children have profile
\begin{equation}
 F_{\mathrm{ch}}(\kappa)
 =\beta\bigl(\kappa-(\tau g-\Delta)\bigr)_+,
 \qquad 0\le\kappa\le g-\Delta.
\label{eq:X-child-profile}
\end{equation}
The bottom interval is filled and has zero cost, and the remaining
interval has length $\ell g$.

Keep the extensively restricted and regularized representative graph
fixed from now on. Pull its event restriction back to its children.
Equation~\eqref{eq:X-frequency} compares normalized representative
and child event weights by subpower factors on every event, so this
pullback is an extensive restriction of the child graph. Uniformize
these children, again extensively, and retain their profile by the
same lemma. This last pruning removes only child events; it does not
alter the already uniform representative graph. Every surviving
child block is therefore still occupied in that fixed representative
graph. Descendant inclusion can be compared within that very block
on each tested gap, regardless of whether different children of the
same representative occupied disjoint branches before pruning.

Each surviving child index stays in one cell $C$ in its chart. Its
surviving times form a subtree of the surviving times of that
representative. On any fixed log gap, choose an occupied block of
this child. Uniformity says that the representative's number of
occupied descendants of that block has its common representative
branching exponent; it is at least the child's number. Therefore,
if $G_{\mathrm{rep}}$ is the representative deficit profile in these
same log units,
\[
 G_{\mathrm{rep}}(t)-G_{\mathrm{rep}}(s)
 \le F_{\mathrm{ch}}(t)-F_{\mathrm{ch}}(s)
 \qquad(0\le s<t\le g-\Delta).
\]
First make this comparison on the finite tested gaps and then use
the Lipschitz profile limits. It implies
$G_{\mathrm{rep}}'\le F_{\mathrm{ch}}'$ almost everywhere. Since
$P$ is increasing, the representative cost is at most
\begin{equation}
 \int_0^{g-\Delta}P(G_{\mathrm{rep}}')
 \le \ell gP(\beta).
\label{eq:X-representative-cost}
\end{equation}
This comparison uses the children on every gap, not just their total
number of events.

Apply the sharp critical bound at resolution $D_0/D$ using
\eqref{eq:X-representative-clustering} and
\eqref{eq:X-representative-cost}. The resulting upper bound for
raw average multiplicity is
\begin{align*}
 &N^{o(1)}
   \bigl[S'D_0^{-B_*}D^d\bigr]
   (D_0/D)^h N^{\ell gP(\beta)}\\
 &\hspace{25mm}=S'D^{d-h}N^{o(1)},
\end{align*}
because $B_*=h+\ell P(\beta)$. This contradicts the strict excess
on the restricted support. Hence
\eqref{eq:X-representative-multiplicity} holds typically. At each fixed
excess margin its exceptional fraction is power-small, by
Convention~\ref{conv:slow-diagonal}. The subpower comparison of the two
event laws therefore makes it negligible for the child events as well.

\paragraph{The packet comparison.}
Fix a surviving packet and one of its filled coarse time bins. Its
mass is at least $S'Z'N^{-o(1)}$. Its child positions lie in a
rectangle of thin width one and long width $Z'$ in $D_0$ cell
units. Replacing a child by its representative moves it by
$O(D)$ such units, by \eqref{eq:X-cell-diameter}. Hence at this
time the packet hashes into at most
\[
 N^{o(1)}O(Z'/D)
\]
spatial cells. In its typical hashed cell its assigned child mass
is therefore at least $S'DN^{-o(1)}$. To make this last typicality
explicit, cells of packet mass below $S'DN^{-\epsilon/4}$ have
total weight at most
$N^{o(1)}(Z'/D)S'DN^{-\epsilon/4}$ at this time, which is a
vanishing fraction of its packet mass. Sum this estimate over
packets and filled bins; assignments are disjoint on each block.

Every child remaining in the test satisfies
\eqref{eq:X-bad-predicate}. Partition the mass of one packet in a
hashed cell according to $C$. Its contribution from $C$ is at most
the original $n(C\cap P)$, hence is at most
$n_CD^{-x}N^{-\epsilon}$. All eligible $C$ representatives are
occupied in that hashed cell. Outside the representative exceptional
set their total $n_C$ is at most
$S'D^{d-h}N^{\epsilon/4}$. The packet mass in that cell is
therefore at most
\[
 S'D^{d-h-x}N^{-3\epsilon/4}
 =S'DN^{-3\epsilon/4},
\]
using $d-h-x=1$. This contradicts the preceding lower bound on a
set of child events of probability tending to one. It proves
\eqref{eq:X-reference-packet}. Choosing the strict margin through
a slow diagonal gives its $o(1)$ form and a vanishing exceptional
event fraction. All selections used only extensive restrictions and
the original upper bounds, so the same proof applies after any
extensive regular refinement of the query law.

\paragraph{From reference mass to normal position.}
The reference-mass claim is now proved. We use full-time packet
containment to place those reference indices near the query in the
finer label's normal-position bins, and then apply cross-entropy.
Return to original coordinates. Let $M$ be the queried matrix and
$\widetilde M$ any reference index in $C\cap P$ counted in
\eqref{eq:X-reference-packet}. Write
$C_j=L_j(c)N^{-a}$ for the four cell widths. Full-time containment
in the packet over $I_b$ gives, at the actual query time $t$ and
in the packet normal $\vartheta_P$,
\[
 |(1,\vartheta_P)\cdot(M-\widetilde M)(t)|
 \le N^{-\rho-wb+o(1)}
 = C_XN^{-\Delta+o(1)}.
\]
This uses containment of each assigned line throughout $I_b$, not
an assumption that each of its shadings meets the query time.

Inside $C$, replacing its frozen frame by an actual $e\in e_c$
changes difference bounds only by dimension constants: the identity
\begin{equation}
X_{t_c+\delta t,\vartheta_c+\delta\vartheta}
 =X_{t_c,\vartheta_c}+\delta t\,U_{t_c,\vartheta_c}
       +\delta\vartheta\,Y_{t_c,\vartheta_c}
       +\delta t\,\delta\vartheta\,V
\label{eq:X-frame-identity}
\end{equation}
has $|\delta t|\lesssim R_c$, $|\delta\vartheta|\lesssim\Theta_c$, and
$C_U=C_X/R_c$, $C_Y=C_X/\Theta_c$,
$C_V=C_X/(R_c\Theta_c)$. Changing $\vartheta_P$ to the query
normal adds at most $\Theta_bN^{o(1)}C_Y$, whose ratio to $C_X$
is $N^{-\tau g+o(1)}$. Changing the query to the center of
$L=e_{c+\sigma}$ adds, relative to $C_X$, at most
\begin{equation}
 N^{-\tau\sigma+o(1)}
 +N^{-\ell\sigma+o(1)}+N^{-\sigma+o(1)}.
\label{eq:X-finer-frame-errors}
\end{equation}
The three terms are the normal, time, and mixed terms of
\eqref{eq:X-frame-identity}. Our choices make
$\tau g,\tau\sigma,\ell\sigma$ strictly greater than
$\Delta$. Consequently all the reference indices counted in
\eqref{eq:X-reference-packet} lie within a subpower number of
additional-precision $X_L$ bins of the queried bin.

\paragraph{Reference cross-entropy.}
For each fixed base $(C,L)$ let $p_{C,L}$ be the histogram obtained
by evaluating $X_L$ on \emph{all} indices in
\eqref{eq:X-reference-mass}, with their weights divided by $n_C$.
Let $q_{C,L}$ be the actual conditional distribution of the query
bin under the event law. These distributions need not agree.
The preceding geometry and \eqref{eq:X-reference-packet} say that,
outside a vanishing event fraction, the $p_{C,L}$ mass in a
subpower enlargement of the queried bin is at least
$D^{-x-o(1)}$.

Choose an integer $r_N=N^{o(1)}$ large enough for that enlargement
on the good events. Replace $p_{C,L}$ by its convolution with the
uniform distribution on $\{-r_N,\ldots,r_N\}$. At a good query
bin the new probability is at least $D^{-x}N^{-o(1)}$.
The image of a full $C$ under $X_L$ has length $O(C_X)$ by
\eqref{eq:X-frame-identity}; hence there are only
$N^{\Delta+o(1)}$ possible bins, including the enlargement.
Take the equal-weight mixture of the convolved histogram and the
uniform law on this finite support, and call it $\widetilde p_{C,L}$.
It is a positive reference probability on every possible query
bin. On good events
\[
 -\log\widetilde p_{C,L}(X_L)
 \le (x\Delta+o(1))\log N,
\]
and on every event it is $O(\log N)$. The vanishing exceptional
fraction thus has vanishing normalized contribution.

Nonnegativity of relative entropy, on each conditional base, gives
\begin{align*}
 \mathrm H\bigl([X_L]_{c,a+\Delta}\mid C,L\bigr)
 &\le
 \E\bigl[-\log\widetilde p_{C,L}
                    ([X_L]_{c,a+\Delta})\bigr]\\
 &\le (x\Delta+o(1))\log N.
\end{align*}
The auxiliary conditioning can now be removed with a conditional
count, rather than the global number of charts. Given $Q(c,a)$ and
$L$, the coarser label $e_c$ and its interval $I_c$ are determined.
In this fixed normalization, the $Q(c,a)$ cell has diameter
$N^{-g+\Delta}$ times its isotropic parent diameter, so it meets
only boundedly many parent cells. The possible bounded grid and
boundary choices have the same property; any further subpower
routing multiplicity contributes only $o(\log N)$. Hence their
conditional entropy given $Q(c,a),L$ is $o(\log N)$, even if the
global collection of parents is polynomial. No $I_b$ or packet-owner
label is conditioned upon in the reference cross-entropy: those
labels were only used to prove the reference-mass inequality, and
the histogram itself uses only $C,L$. The chain rule therefore gives
\eqref{eq:conditional-X-bound} with exactly its stated conditioning.
No independence of the matrix and the finer event label has been used.
\end{proof}

Both lemmas are formulated with fixed positive gaps. Their constants
and strict margins can first be imposed on any finite collection of
scale and offset tests. The cleaning procedure and the stationary
construction then permit countable diagonals, with errors tending to
zero before a tested gap is shrunk. In particular
Lemma~\ref{lem:conditional-X} supplies an upper bound for a
conditional entropy increment; it does not identify that increment
with a Frostman exponent without the additional joint-cell
regularizations used below.
